\documentclass[
  bibliography=totoc,     % Literatur im Inhaltsverzeichnis
  captions=tableheading,  % Tabellenüberschriften
  titlepage=firstiscover, % Titelseite ist Deckblatt
]{scrartcl}

%irgendetwas mit Tabelen und Figuren anders Nummerieren 
\usepackage{chngcntr}
\usepackage{longtable} 

\usepackage{titling}

%Textdatein einfügen
\usepackage{verbatim}

% Paket float verbessern
\usepackage{scrhack}

% Warnung, falls nochmal kompiliert werden muss
\usepackage[aux]{rerunfilecheck}

% unverzichtbare Mathe-Befehle
\usepackage{amsmath}
% viele Mathe-Symbole
\usepackage{amssymb}
% Erweiterungen für amsmath
\usepackage{mathtools}

% Fonteinstellungen
\usepackage{fontspec}
% Latin Modern Fonts werden automatisch geladen
% Alternativ zum Beispiel:
%\setromanfont{Libertinus Serif}
%\setsansfont{Libertinus Sans}
%\setmonofont{Libertinus Mono}

% Wenn man andere Schriftarten gesetzt hat,
% sollte man das Seiten-Layout neu berechnen lassen
\recalctypearea{}

% deutsche Spracheinstellungen
\usepackage[ngerman]{babel}


\usepackage[
  math-style=ISO,    % ┐
  bold-style=ISO,    % │
  sans-style=italic, % │ ISO-Standard folgen
  nabla=upright,     % │
  partial=upright,   % ┘
  warnings-off={           % ┐
    mathtools-colon,       % │ unnötige Warnungen ausschalten
    mathtools-overbracket, % │
  },                       % ┘
]{unicode-math}

% traditionelle Fonts für Mathematik
\setmathfont{Latin Modern Math}
% Alternativ zum Beispiel:
%\setmathfont{Libertinus Math}

\setmathfont{XITS Math}[range={scr, bfscr}]
\setmathfont{XITS Math}[range={cal, bfcal}, StylisticSet=1]

% Zahlen und Einheiten
\usepackage[
  locale=DE,                   % deutsche Einstellungen
  separate-uncertainty=true,   % immer Unsicherheit mit \pm
  per-mode=symbol-or-fraction, % / in inline math, fraction in display math
]{siunitx}

\DeclareSIUnit{\channel}{Channel}
\DeclareSIUnit{\year}{a}

% chemische Formeln
\usepackage[
  version=4,
  math-greek=default, % ┐ mit unicode-math zusammenarbeiten
  text-greek=default, % ┘
]{mhchem}

% richtige Anführungszeichen
\usepackage[autostyle]{csquotes}

% schöne Brüche im Text
\usepackage{xfrac}

% Standardplatzierung für Floats einstellen
\usepackage{float}
\floatplacement{figure}{htbp}
\floatplacement{table}{htbp}

% Floats innerhalb einer Section halten
\usepackage[
  section, % Floats innerhalb der Section halten
  below,   % unterhalb der Section aber auf der selben Seite ist ok
]{placeins}

% Seite drehen für breite Tabellen: landscape Umgebung
\usepackage{pdflscape}

% Captions schöner machen.
\usepackage[
  labelfont=bf,        % Tabelle x: Abbildung y: ist jetzt fett
  font=small,          % Schrift etwas kleiner als Dokument
  width=0.9\textwidth, % maximale Breite einer Caption schmaler
]{caption}
% subfigure, subtable, subref
\usepackage{subcaption}

% Grafiken können eingebunden werden
\usepackage{graphicx}
\usepackage{wrapfig}

% schöne Tabellen
\usepackage{booktabs}
\usepackage[table]{xcolor}

% Verbesserungen am Schriftbild
\usepackage{microtype}

% Literaturverzeichnis
\usepackage[
  backend=biber,
  sorting=none
]{biblatex}
% Quellendatenbank
\addbibresource{lit.bib}
\addbibresource{programme.bib}

% Hyperlinks im Dokument
\usepackage[
  german,
  unicode,        % Unicode in PDF-Attributen erlauben
  pdfusetitle,    % Titel, Autoren und Datum als PDF-Attribute
  pdfcreator={},  % ┐ PDF-Attribute säubern
  pdfproducer={}, % ┘
]{hyperref}
% erweiterte Bookmarks im PDF
\usepackage{bookmark}

% Trennung von Wörtern mit Strichen
\usepackage[shortcuts]{extdash}

%\setcounter{tocdepth}{3} % + subsubsections



\author{%
  Benedikt Lütke Lanfer \\%
  \href{mailto:benedikt.luetkelanfer@tu-dortmund.de}{benedikt.luetkelanfer@tu-dortmund.de}%
  \and%
  Enno Wellmann \\%
  \href{mailto:enno.wellmann@tu-dortmund.de}{enno.wellmann@tu-dortmund.de}%
}
\publishers{TU Dortmund – Fakultät Physik}


\newcommand*\diff{\mathop{}\!\mathrm{d}}

\NewDocumentCommand \OverfullCenter {+m} {
\noindent\makebox[\linewidth]{#1} }

\usepackage{adjustbox}


% %Tabellen und Figuren Einstellung
% \counterwithout{table}{section}
% \counterwithout{figure}{section}
% \renewcommand{\thetable}{\Roman{table}}
% \renewcommand{\thefigure}{\Roman{figure}}

%Richtiges Einrücken
\setlength{\parindent}{0pt}


\title{V46:\\ Faraday-Effekt an Halbleitern}
\author{Benedikt Lütke Lanfer \and Enno Wellmann}
\date{10. Juni 2024}
\publishers{TU Dortmund – Fakultät Physik}

\begin{document}
\begin{titlingpage}
    \begin{center}
        \begin{Huge}
            \textbf{\thetitle\\}
        \end{Huge}
    \end{center}
    \vspace{4cm}
    \includegraphics[width=\textwidth]{Bilder/Logo_TU.png} \\
    \vspace{4cm}
    \begin{center}
        \begin{huge}
            \theauthor\\
        \end{huge}
        \vspace{0.5cm}
        \begin{Large}
            benedikt.luetkelanfer@tu-dortmund.de\\
            enno.wellmann@tu-dortmund.de\\
            \vspace{1.4cm}
            Bearbeitet: \today\\
            Durchgeführt: \thedate\\
            TU Dortmund – Fakultät Physik\\
        \end{Large}
    \end{center}
\end{titlingpage}
\tableofcontents
\newpage

\section{Zielsetzung}

%---------------------------------------------------------------------------------------------------------------------------------------------------------------%

\section{Theorie}

\subsection{Fehlerrechnung}
Für die Fehlerrechnung werden alle \textbf{Mittelwerte} von $N$ Messungen folgendermaßen berechnet:

\begin{equation}
    \overline{x} = \frac{1}{N} \cdot \sum_{i=1}^N x_i
    \label{eqn:Mittelwert}
\end{equation}

und alle \textbf{Standardabweichungen zum Mittelwert} mit:

\begin{equation}
    \increment\overline{x} = \sqrt{\frac{1}{N\cdot(N-1)}\cdot\sum_{i=1}^N (x_i-\overline{x})^2}
    \label{eqn:St_Mittelwert}
\end{equation}

Der Fehler für zusammenhängende Messwerte wird dann mit der \textbf{Gaußschen Fehlerfortpflanzung} berechnet:

\begin{equation}
    \increment{f} = \sqrt{ \sum_{i = 1}^{N}  \biggl(\frac{\partial{f}}{\partial{x_i}}\biggr)^2\cdot(\increment{x_i})^2}
    \label{eqn:Gauss}
\end{equation}

Die Fehlerfortpflanzung wird mit Uncertainties in Python \cite{uncertainties} ermittelt.

%---------------------------------------------------------------------------------------------------------------------------------------------------------------%

\section{Durchführung}


%---------------------------------------------------------------------------------------------------------------------------------------------------------------%
\newpage
\section{Auswertung}
\subsection{Bestimmung des B-Feldes}
Die Messung des B-Feldes mittels einer Hall-Sonde ergibt die in Tabelle \eqref{tab:B_Feld} folgende Werte:

\begin{table}[H]
	\centering
	\begin{tabular}{c c}
		\toprule
		$d \, [\unit{\milli\meter}]$ & $B \, [\unit{\milli\tesla}] $  \\
		\midrule
        170 & 1 \\
        160 & 0 \\
        150 & 0 \\
        140 & 0 \\
        130 & 1 \\
        120 & 6 \\
        112 & 92 \\
        110 & 176 \\
        108 & 280 \\
        105 & 382 \\
        102 & 418 \\
        100 & 426 \\
        98  & 428 \\
        95  & 420 \\
        90  & 340 \\
        87  & 236 \\
        85  & 97 \\
        80  & 17 \\
        70  & 0  \\
		\bottomrule
	\end{tabular}
    \label{tab:B_Feld}
\end{table}

Diese Werte werden einmal geplottet \eqref{fig:B_Feld}, welches eine gut sichtbare Parabel förmige Struktur ergibt. 
Jedoch wurden die Werte von $\qty{140}{\milli\meter}$ bis $\qty{170}{\milli\meter}$ ausgelassen das diese Vernachlässigt werden können. 
Das maximale B-Feld ergibt sich dabei mit $\qty{428}{\milli\tesla}$ bei $\qty{98}{\milli\meter}$. 

\begin{figure}[H]
	\centering
	\includegraphics[width=\textwidth]{build/B_Feld.pdf}
	\caption{Messdaten B-Feld}\label{fig:B_Feld}
\end{figure}

\subsection{Bestimmung der effektiven Masse}
Die Differenz der einzelnen aufgenommen Messwerte der Winkel je nach B-Feld Ausrichtung werden zuerst mittels 

\begin{equation}
    \theta=\frac{\theta_2 - \theta_1}{2L} 
\end{equation}

auf die Länge $L$ normiert. 
Der Faktor $2$ kommt daher, dass die Differenz des B-Feldes doppelt so stark ist wie das vorher bestimmte maximale B-Feld ist. 

\begin{table}[H]
	\centering
	\caption{Ergebnisse Vollenergienachweiswahrscheinlichkeit}
	\begin{tabular}{|c c | c c | c c|}
		\toprule
		\multicolumn{2}{|c|}{$GaAs_n$}  & \multicolumn{2}{|c|}{$GaAs_r \, $}  & \multicolumn{2}{|c|}{$GaAs_n \, $}   \\
        \multicolumn{2}{|c|}{$N=\qty[per-mode=fraction]{2.8e18}{\per\cubic\centi\meter} $}  & \multicolumn{2}{|c|}{$ $}  & \multicolumn{2}{|c|}{$N=\qty[per-mode=fraction]{1.28e18}{\per\cubic\centi\meter} $}   \\
        \midrule
        $\theta_1$ & $\theta_2 $ & $\theta_1 $ & $ \theta_2 $ & $\theta_1 $ & $ \theta_2 $ \\
		\midrule
            \num{77.50} & \num{88.60} & \num{70.50} & \num{93.50} & \num{78.42} & \num{87.05} \\
            \num{78.00} & \num{86.12} & \num{75.08} & \num{90.42} & \num{80.17} & \num{86.93} \\
            \num{74.17} & \num{87.63} & \num{78.50} & \num{90.50} & \num{80.25} & \num{88.08} \\
            \num{77.08} & \num{86.48} & \num{78.53} & \num{87.00} & \num{80.00} & \num{86.75} \\
            \num{71.33} & \num{87.27} & \num{74.83} & \num{82.50} & \num{75.33} & \num{81.08} \\
            \num{69.50} & \num{82.32} & \num{74.67} & \num{79.17} & \num{72.78} & \num{78.83} \\
            \num{44.08} & \num{51.50} & \num{48.50} & \num{52.50} & \num{45.25} & \num{50.25} \\
            \num{27.50} & \num{40.41} & \num{29.17} & \num{32.33} & \num{28.25} & \num{36.00} \\
            \num{63.75} & \num{78.22} & \num{67.83} & \num{73.17} & \num{66.60} & \num{73.08} \\
		\bottomrule
	\end{tabular}
	\label{tab:data2}
\end{table}

\begin{figure}[H]
	\centering
	\includegraphics[width=\textwidth]{build/Messdaten.pdf}
	\caption{Messdaten Proben}\label{fig:Aufbau}
\end{figure}

\begin{figure}[H]
	\centering
	\includegraphics[width=\textwidth]{build/plot1.pdf}
	\caption{Ausgleichsgerade}\label{fig:Aufbau}
\end{figure}

\begin{figure}[H]
	\centering
	\includegraphics[width=\textwidth]{build/plot2.pdf}
	\caption{Ausgleichsgerade}\label{fig:Aufbau}
\end{figure}

%---------------------------------------------------------------------------------------------------------------------------------------------------------------%

\section{Diskussion}


%---------------------------------------------------------------------------------------------------------------------------------------------------------------%
\newpage
\printbibliography

\end{document}